\documentclass[aspectratio=169]{beamer}
\usetheme{Metropolis}
\setbeamertemplate{blocks}[default]

\usepackage[utf8]{inputenc}
\usepackage[spanish,es-tabla]{babel}
\usepackage{lmodern}
\usepackage{amsmath,amssymb}
\usepackage{booktabs}
\usepackage{array}
\usepackage{graphicx}
\usepackage{tikz}
\usetikzlibrary{arrows.meta,positioning,fit,backgrounds}
\usepackage{xcolor}
\usepackage{colortbl}
\usepackage{listings}

\definecolor{entra}{HTML}{2E7D32}
\definecolor{sale}{HTML}{8C1D40}
\definecolor{pivcolor}{HTML}{D2691E}
\definecolor{basecolor}{HTML}{3B7EA1}
\definecolor{podado}{HTML}{B0413E}

\lstset{
  basicstyle=\ttfamily\scriptsize,
  breaklines=true,
  columns=fullflexible,
  keywordstyle=\color{basecolor}\bfseries,
  commentstyle=\color{gray},
  stringstyle=\color{entra}
}

\title{Complemento a la Clase 4: Infactibilidad y el Método de Dos Fases}
\subtitle{El cuarto caso especial de la Programación Lineal}
\author{Métodos de Optimización --- Maestría}
\institute{Universidad Pontificia Bolivariana}
\date{}

\begin{document}

\maketitle

% ================= SECCION 1 =================
\section{El cuarto caso especial}

\begin{frame}{Lo que ya vimos en la Clase 4}
En la Clase 4 cubrimos tres desenlaces especiales del simplex:

\begin{enumerate}
\item \textbf{Óptimos múltiples}: un coeficiente exactamente cero en la fila $z$ de una no básica al terminar.
\item \textbf{No acotado}: una columna entrante sin ninguna razón positiva --- $z\to\infty$.
\item \textbf{Degeneración}: empate en la razón mínima, riesgo de ciclado, regla de Bland.
\end{enumerate}

\begin{alertblock}{Lo que faltó}
Un cuarto caso: \textbf{infactibilidad} --- la región factible está \emph{vacía}, no existe ningún punto que cumpla todas las restricciones a la vez. No es un ``tipo'' de degeneración ni de ningún otro caso ya visto: es una falla distinta y anterior.
\end{alertblock}
\end{frame}

\begin{frame}{¿Por qué nunca lo vimos en la Clase 4?}
No fue casualidad --- fue consecuencia directa de cómo construimos \emph{todos} los ejemplos de esa clase.

\begin{block}{El truco que usamos siempre}
Con restricciones $\leq$ y $b_i\geq0$, la variable de holgura da un arranque \textbf{garantizado}: $x=0$, todas las holguras $=b_i\geq0$. El sistema ya viene diagonalizado, factible, sin ningún trabajo.
\end{block}

Nunca introdujimos una restricción $\geq$ o $=$ --- así que nunca le dimos al simplex la oportunidad de toparse con un problema sin solución. Ese arranque garantizado es, precisamente, lo que se pierde con $\geq$ o $=$.
\end{frame}

% ================= SECCION 2 =================
\section{El problema de arrancar con restricciones $\geq$}

\begin{frame}{El excedente no resuelve el problema de arranque}
Para $x_1\geq3$, la forma estándar es $x_1-s_1=3$ con $s_1\geq0$ (variable de \textbf{exceso}, no de holgura). Algebraicamente es perfecta. El problema es \textbf{arrancar} el simplex desde ahí.

Intenta el mismo truco de siempre: fija $x_1=0$ (no básica), trata $s_1$ como básica:
\[
-s_1=3 \quad\Rightarrow\quad s_1=-3
\]

\begin{alertblock}{Viola $s_1\geq0$}
Ese punto de arranque no es factible. El signo del excedente ($-s_1$, no $+s_1$) es justo la diferencia con la holgura de $\leq$.
\end{alertblock}

\textbf{¿Y usar $x_1$ mismo como básica?} Si $x_1\geq3$ fuera la \emph{única} restricción, sí --- pero en un problema real $x_1$ suele aparecer en varias filas más, y una variable básica necesita una columna ``limpia'' (un $1$ en su fila, $0$ en las demás). La columna real de $x_1$ casi nunca cumple eso.
\end{frame}

\begin{frame}{La variable artificial: un truco puramente mecánico}
Se agrega una \textbf{variable artificial} $a_1\geq0$, construida a propósito para tener columna limpia:
\[
x_1-s_1+a_1=3
\]

Fija $x_1=s_1=0$ (no básicas), lee $a_1=3\geq0$ directo --- \textbf{factible de inmediato}, exactamente el mismo arranque gratis que dábamos por sentado con las holguras.

\begin{block}{El precio}
$a_1$ no es una cantidad real del problema --- es un valor ``de mentiras'' que presta un punto de partida. Por eso hay que \textbf{forzarla a valer cero} antes de confiar en cualquier resultado. Eso es exactamente lo que hace la \textbf{Fase 1}.
\end{block}
\end{frame}

% ================= SECCION 3 =================
\section{El método de Dos Fases}

\begin{frame}{Fase 1: minimizar la suma de artificiales}
\begin{block}{Fase 1}
\small Se resuelve un simplex auxiliar cuyo único objetivo es
\[
\min\ w=\sum_i a_i \qquad\text{(equivalentemente, } \max\ w'=-w\text{)}
\]
usando \textbf{el mismo simplex de siempre} --- misma regla de entrada, misma razón mínima, mismo Gauss-Jordan.
\normalsize
\end{block}

\vspace{-0.1cm}
\begin{alertblock}{El criterio que decide todo}
\small Al terminar Fase 1:
\begin{itemize}
\item Si $w=0$ (todas las artificiales en cero): se encontró un punto \textbf{genuinamente factible}. Se descartan las columnas artificiales y se pasa a Fase 2.
\item Si $w>0$ (el mínimo posible de $w$ es positivo): \textbf{no existe} ningún punto que cumpla todas las restricciones. El problema es \textbf{infactible}. Punto final --- no hay Fase 2.
\end{itemize}
\normalsize
\end{alertblock}
\end{frame}

\begin{frame}{Fase 2: optimizar el objetivo real}
Si Fase 1 terminó con $w=0$:

\begin{enumerate}
\item Se \textbf{descartan las columnas de las variables artificiales} por completo (nunca vuelven a entrar a la base).
\item Se \textbf{restaura la fila $z$ real} del problema original, sobre la base factible que dejó la Fase 1.
\item Se continúa el simplex \textbf{normal} (el de la Clase 4) hasta el criterio de optimalidad de siempre.
\end{enumerate}

\begin{alertblock}{La idea completa, en una frase}
Dos Fases no es un algoritmo nuevo: es el simplex de la Clase 4, corrido \textbf{dos veces} --- primero con un objetivo prestado que solo busca un punto de partida factible, y luego con el objetivo real, sobre la base que la primera corrida dejó lista.
\end{alertblock}
\end{frame}

% ================= SECCION 4 =================
\section{Ejemplo A: factible tras Fase 1}

\begin{frame}{Ejemplo A: planteamiento}
\begin{block}{Problema}
\[
\max\ z=2x_1+3x_2 \quad\text{s.a.}\quad x_1+x_2\leq4,\ \ x_1+2x_2\geq4,\ \ x_1,x_2\geq0
\]
\end{block}

En el origen ($x_1=x_2=0$), la segunda restricción da $0\geq4$: \textbf{falsa}. El origen no es factible --- por eso hace falta Fase 1 aquí, a diferencia de todos los ejemplos de la Clase 4.

Forma estándar: $x_1+x_2+s_1=4$ (holgura normal) y $x_1+2x_2-s_2+a_1=4$ (excedente + artificial).
\end{frame}

\begin{frame}{Ejemplo A: Fase 1, tableau inicial}
\begin{center}
\renewcommand{\arraystretch}{1.25}
\begin{tabular}{c|cccc|c}
\textbf{Base} & $x_1$ & $x_2$ & $s_1$ & $s_2$ & $b$ \\
\hline
$s_1$ & $1$ & $1$ & $1$ & $0$ & $4$ \\
$a_1$ & $1$ & $2$ & $0$ & $-1$ & $4$ \\
\hline
$w'$ & $-1$ & \cellcolor{entra!20}$-2$ & $0$ & $1$ & $-4$ \\
\end{tabular}
\end{center}

\vspace{0.2cm}
(Columna $a_1$ omitida tras eliminarla de su propia fila en la fila $w'$; su coeficiente ya es $0$ ahí por construcción.) Coeficiente más negativo: $x_2$ ($-2$) $\to$ entra $x_2$. Razón mínima: $\min(4/1,\ 4/2)=2\to$ sale $a_1$.
\end{frame}

\begin{frame}{Ejemplo A: Fase 1, tras pivotear --- ya factible}
\begin{center}
\renewcommand{\arraystretch}{1.25}
\begin{tabular}{c|cccc|c}
\textbf{Base} & $x_1$ & $x_2$ & $s_1$ & $s_2$ & $b$ \\
\hline
$s_1$ & $1/2$ & $0$ & $1$ & $1/2$ & $2$ \\
$x_2$ & $1/2$ & $1$ & $0$ & $-1/2$ & $2$ \\
\hline
$w'$ & $0$ & $0$ & $0$ & $0$ & $\mathbf{0}$ \\
\end{tabular}
\end{center}

\begin{alertblock}{$w=0$: FACTIBLE}
Se encontró un punto real: $x_1=0,\ x_2=2$ ($a_1$ ya no aparece, quedó en $0$). Sin coeficientes negativos en $w'$: Fase 1 terminada. \textbf{Se descarta la columna $a_1$} y se pasa a Fase 2 con la base $\{s_1,x_2\}$ tal cual quedó.
\end{alertblock}
\end{frame}

\begin{frame}{Ejemplo A: Fase 2, tableau inicial}
Se restaura $z=2x_1+3x_2$ sobre la base $\{s_1,x_2\}$ heredada de Fase 1 (columna $a_1$ ya no existe):

\begin{center}
\renewcommand{\arraystretch}{1.25}
\begin{tabular}{c|ccc|c}
\textbf{Base} & $x_1$ & $x_2$ & $s_2$ & $b$ \\
\hline
$s_1$ & $1/2$ & $0$ & $1/2$ & $2$ \\
$x_2$ & $1/2$ & $1$ & $-1/2$ & $2$ \\
\hline
$z$ & \cellcolor{entra!20}$-1/2$ & $0$ & $-3/2$ & $6$ \\
\end{tabular}
\end{center}

$z=6+0.5x_1+1.5s_2$: todavía no óptimo (dos coeficientes positivos en la lectura real). Entra $s_2$ (mayor). Razón mínima: solo $s_1$ tiene coeficiente positivo en $s_2$ $\to$ $\min(2/0.5)=4\to$ sale $s_1$.
\end{frame}

\begin{frame}{Ejemplo A: Fase 2, óptimo}
\begin{center}
\renewcommand{\arraystretch}{1.25}
\begin{tabular}{c|ccc|c}
\textbf{Base} & $x_1$ & $x_2$ & $s_2$ & $b$ \\
\hline
$s_2$ & $1$ & $0$ & $1$ & $4$ \\
$x_2$ & $1$ & $1$ & $0$ & $4$ \\
\hline
$z$ & $1$ & $0$ & $0$ & $\mathbf{12}$ \\
\end{tabular}
\end{center}

Sin coeficientes positivos reales (fila $z$: $z=12-x_1$): \textbf{óptimo}. $x_1^\ast=0,\ x_2^\ast=4,\ z^\ast=12$.

\begin{block}{Verificación geométrica}
En $(0,4)$: $x_1+x_2=4\leq4$ (holgura $s_1=0$, se agotó) y $x_1+2x_2=8\geq4$ (excedente $s_2=4$, sobra). Ambas restricciones se cumplen --- consistente con lo que el tableau final reporta.
\end{block}
\end{frame}

% ================= SECCION 5 =================
\section{Ejemplo B: infactible}

\begin{frame}{Ejemplo B: el nodo del árbol de la Clase 5}
Recordemos el nodo que generó tu pregunta original: dentro del árbol de Branch and Bound del Ejemplo 2 de la Clase 5, un nodo heredaba
\[
2x_1+3x_2\leq12 \quad(\text{original}), \qquad x_1\geq2, \qquad x_2\geq3
\]
(más $x_1\leq2$, que aquí tratamos como cota de variable, no como fila --- igual que discutimos en el chat).

Forma estándar de las partes ``difíciles'': $x_1-s_2+a_1=2$ y $x_2-s_3+a_2=3$, cada una con su propia artificial.
\end{frame}

\begin{frame}{Ejemplo B: Fase 1, tableau inicial}
\begin{center}
\renewcommand{\arraystretch}{1.2}
\small
\begin{tabular}{c|ccccc|c}
\textbf{Base} & $x_1$ & $x_2$ & $s_1$ & $s_2$ & $s_3$ & $b$ \\
\hline
$s_1$ & $2$ & $3$ & $1$ & $0$ & $0$ & $12$ \\
$a_1$ & $1$ & $0$ & $0$ & $-1$ & $0$ & $2$ \\
$a_2$ & $0$ & $1$ & $0$ & $0$ & $-1$ & $3$ \\
\hline
$w'$ & \cellcolor{entra!20}$-1$ & $-1$ & $0$ & $1$ & $1$ & $-5$ \\
\end{tabular}
\end{center}
\normalsize

\vspace{0.2cm}
Dos candidatas empatadas ($x_1,x_2$ ambas con $-1$); tomamos $x_1$. Razón mínima: $\min(12/2,\ 2/1)=2\to$ sale $a_1$.
\end{frame}

\begin{frame}{Ejemplo B: Fase 1, iteración 1}
\begin{center}
\renewcommand{\arraystretch}{1.2}
\small
\begin{tabular}{c|ccccc|c}
\textbf{Base} & $x_1$ & $x_2$ & $s_1$ & $s_2$ & $s_3$ & $b$ \\
\hline
$s_1$ & $0$ & $3$ & $1$ & $2$ & $0$ & $8$ \\
$x_1$ & $1$ & $0$ & $0$ & $-1$ & $0$ & $2$ \\
$a_2$ & $0$ & $1$ & $0$ & $0$ & $-1$ & $3$ \\
\hline
$w'$ & $0$ & \cellcolor{entra!20}$-1$ & $0$ & $0$ & $1$ & $-3$ \\
\end{tabular}
\end{center}
\normalsize

\vspace{0.2cm}
Entra $x_2$ (único coeficiente negativo). Razón mínima: $\min(8/3,\ 3/1)=8/3\to$ sale $s_1$.
\end{frame}

\begin{frame}{Ejemplo B: Fase 1, óptimo --- $w>0$}
\begin{center}
\renewcommand{\arraystretch}{1.2}
\small
\begin{tabular}{c|ccccc|c}
\textbf{Base} & $x_1$ & $x_2$ & $s_1$ & $s_2$ & $s_3$ & $b$ \\
\hline
$x_2$ & $0$ & $1$ & $1/3$ & $2/3$ & $0$ & $8/3$ \\
$x_1$ & $1$ & $0$ & $0$ & $-1$ & $0$ & $2$ \\
$a_2$ & $0$ & $0$ & $-1/3$ & $-2/3$ & $-1$ & $\mathbf{1/3}$ \\
\hline
$w'$ & $0$ & $0$ & $1/3$ & $2/3$ & $1$ & $\mathbf{-1/3}$ \\
\end{tabular}
\end{center}
\normalsize

\begin{alertblock}{$w=1/3>0$: INFACTIBLE}
Sin coeficientes negativos en $w'$: Fase 1 terminó, pero el mínimo posible de $a_1+a_2$ es $1/3$, \textbf{no cero}. No existe ningún punto que cumpla las tres restricciones simultáneamente --- esto explica por qué scipy reportaba ese nodo como infeasible.
\end{alertblock}
\end{frame}

\begin{frame}{Ejemplo B: leyendo la infactibilidad en números reales}
En el óptimo de Fase 1: $x_1=2,\ x_2=8/3\approx2.667,\ a_2=1/3$.

\begin{block}{Interpretación}
El simplex \textbf{sí encontró} el mejor compromiso posible: con $x_1=2$ fijo, la restricción de recursos $2x_1+3x_2\leq12$ solo permite $x_2\leq8/3$. Pero la restricción heredada exige $x_2\geq3$. Le faltan exactamente $3-8/3=1/3$ unidades para cumplirla --- eso es literalmente lo que mide $a_2=1/3$: \textbf{cuánto le falta} a la mejor solución posible para ser genuinamente factible.
\end{block}

Esto coincide exactamente con lo que ya habíamos deducido a mano en el chat, sin Fase 1: sustituyendo $x_1=2$ en $2x_1+3x_2\leq12$ se obtiene $x_2\leq8/3$, que contradice $x_2\geq3$.
\end{frame}

\begin{frame}{El certificado de infactibilidad}
\begin{block}{No es una conjetura --- es una prueba}
Multiplica la restricción de recursos por $1$ y la restricción $x_2\geq3$ (escrita como $-x_2\leq-3$) por $3$, y súmalas:
\[
1\times(2x_1+3x_2\leq12) \ + \ 3\times(-x_2\leq-3) \quad\text{con } x_1=2 \text{ fijo}
\]
\[
\Rightarrow\quad 3x_2\leq8 \qquad\text{y}\qquad 3x_2\geq9 \qquad\Rightarrow\qquad 9\leq8
\]
\end{block}

\textbf{$9\leq8$ es una contradicción directa.} Esta combinación (con coeficientes $1$ y $3$) es exactamente lo que los coeficientes finales de $s_1$ y $s_3$ en la fila $w'$ ($1/3$ y $1$, escalados) codifican algebraicamente. Se llama \textbf{certificado de infactibilidad} (Lema de Farkas): una combinación de tus propias restricciones que prueba, sin lugar a dudas, que ninguna solución puede existir.
\end{frame}

% ================= SECCION 6 =================
\section{Alternativa: el método de la Gran M}

\begin{frame}{Big M: la misma idea, un solo objetivo}
\begin{block}{Método de la Gran M}
\small En vez de dos fases separadas, se penaliza cada artificial directamente en el objetivo \textbf{real} con un coeficiente enorme $M$:
\[
\max\ z = 2x_1+3x_2 - M a_1 \qquad (M \text{ muy grande, p.ej. } M=1000)
\]
\normalsize
\end{block}

\vspace{-0.1cm}
\small Como $M$ es enorme, el simplex \textbf{siempre} preferirá sacar cualquier artificial de la base antes que tolerar su penalización --- el efecto práctico es casi idéntico a Fase 1, pero en una sola corrida.

\begin{alertblock}{La desventaja}
\small $M$ debe ser ``suficientemente grande'', pero no hay una regla universal de cuánto --- un $M$ mal elegido puede causar errores numéricos (redondeo). Por eso los solvers profesionales casi siempre prefieren Dos Fases o variantes más robustas.
\normalsize
\end{alertblock}
\end{frame}

% ================= SECCION 7 =================
\section{Verificación y síntesis}

\begin{frame}[fragile]{Verificación: \texttt{scipy} detecta la infactibilidad directo}
\begin{lstlisting}[language=Python]
from scipy.optimize import linprog

# Ejemplo B: 2x1+3x2<=12, x1>=2, x2>=3
res = linprog(
    c=[0, 0],
    A_ub=[[2, 3], [-1, 0], [0, -1]],
    b_ub=[12, -2, -3],
    bounds=[(0, None), (0, None)],
    method="highs",
)
print(res.status, res.message)
# -> 2  The problem is infeasible.
\end{lstlisting}

Por dentro, HiGHS corre su propia variante de Fase 1 (o un método de punto interior con detección de infactibilidad) --- el \texttt{status} que devuelve es la misma conclusión que llegamos a mano.
\end{frame}

\begin{frame}{Los cuatro casos especiales de PL --- síntesis actualizada}
\begin{center}
\small
\renewcommand{\arraystretch}{1.35}
\begin{tabular}{@{}p{2.6cm}p{4.6cm}p{4.6cm}@{}}
\toprule
\textbf{Caso} & \textbf{Señal en el tableau} & \textbf{Qué significa}\\
\midrule
Óptimos múltiples & Coeficiente cero en $z$ de una no básica, en el óptimo & Toda una cara/segmento es óptima, no solo un vértice\\
No acotado & Columna entrante sin ninguna razón positiva & $z\to\infty$; falta una restricción en el modelo\\
Degenerado & Empate en la razón mínima & Vértice donde sobran restricciones activas; riesgo de ciclado\\
\textbf{Infactible} & \textbf{Fase 1 termina con $w>0$} & \textbf{La región factible está vacía; ninguna solución existe}\\
\bottomrule
\end{tabular}
\end{center}
\normalsize
\end{frame}

\begin{frame}{Ejercicios propuestos}
\small
\begin{enumerate}
\setlength{\itemsep}{2pt}
\item Resuelva a mano, con Dos Fases, $\max z=x_1+2x_2$ s.a. $x_1+x_2\leq6,\ x_1+3x_2\geq9,\ x_1,x_2\geq0$. Verifique con \texttt{scipy}.
\item Demuestre que $\max z=x_1$ s.a. $x_1\geq5,\ x_1\leq2,\ x_1\geq0$ es infactible, tanto por Fase 1 como con un certificado algebraico directo (sin variables artificiales, en una sola línea).
\item Para el Ejemplo A de esta clase, ¿qué habría pasado en Fase 1 si en vez de elegir $x_2$ como entrante hubiéramos elegido $x_1$ (también negativo)? Verifique si se llega igual a $w=0$.
\item Investigue por qué el método de la Gran M puede fallar numéricamente con una elección de $M$ demasiado grande, y qué alternativa usan los solvers modernos (pista: buscar ``Big M numerical issues linear programming'').
\item Retome el Ejemplo 1 (sin poda genuina) de la Clase 5 y verifique si alguno de sus nodos requería Fase 1. ¿Por qué sí o por qué no?
\item Escriba el certificado de infactibilidad (combinación lineal de restricciones) para el ejercicio 2, siguiendo el mismo estilo del certificado mostrado en esta clase para el Ejemplo B.
\end{enumerate}
\end{frame}

\begin{frame}{Referencias}
\small
\begin{itemize}
\setlength{\itemsep}{1pt}
\item Bertsimas, D., \& Tsitsiklis, J. N. (1997). \emph{Introduction to Linear Optimization}, Cap. 3 (Fase 1, Gran M). Athena Scientific.
\item Chvátal, V. (1983). \emph{Linear Programming}, Cap. 5. W. H. Freeman.
\item Farkas, J. (1902). ``Theorie der einfachen Ungleichungen''. \emph{Journal für die reine und angewandte Mathematik}, 124, 1--27. (Origen del lema de Farkas.)
\item Vanderbei, R. J. (2020). \emph{Linear Programming: Foundations and Extensions} (5th ed.), Cap. 4. Springer.
\item Dantzig, G. B., \& Thapa, M. N. (1997). \emph{Linear Programming 1: Introduction}. Springer.
\item Nocedal, J., \& Wright, S. J. (2006). \emph{Numerical Optimization} (2nd ed.), Cap. 13. Springer.
\end{itemize}
\end{frame}

\begin{frame}{}
\centering
{\Large\textbf{Gracias}}

\vspace{0.4cm}
\normalsize Preguntas y discusión
\end{frame}

\end{document}
